% Options for packages loaded elsewhere
\PassOptionsToPackage{unicode}{hyperref}
\PassOptionsToPackage{hyphens}{url}
\documentclass[
  11pt,
]{article}
\usepackage{xcolor}
\usepackage[left=2.5cm, right=2.5cm, top=2.5cm, bottom=2.5cm]{geometry}
\usepackage{amsmath,amssymb}
\setcounter{secnumdepth}{5}
\usepackage{iftex}
\ifPDFTeX
  \usepackage[T1]{fontenc}
  \usepackage[utf8]{inputenc}
  \usepackage{textcomp} % provide euro and other symbols
\else % if luatex or xetex
  \usepackage{unicode-math} % this also loads fontspec
  \defaultfontfeatures{Scale=MatchLowercase}
  \defaultfontfeatures[\rmfamily]{Ligatures=TeX,Scale=1}
\fi
\usepackage{lmodern}
\ifPDFTeX\else
  % xetex/luatex font selection
\fi
% Use upquote if available, for straight quotes in verbatim environments
\IfFileExists{upquote.sty}{\usepackage{upquote}}{}
\IfFileExists{microtype.sty}{% use microtype if available
  \usepackage[]{microtype}
  \UseMicrotypeSet[protrusion]{basicmath} % disable protrusion for tt fonts
}{}
\makeatletter
\@ifundefined{KOMAClassName}{% if non-KOMA class
  \IfFileExists{parskip.sty}{%
    \usepackage{parskip}
  }{% else
    \setlength{\parindent}{0pt}
    \setlength{\parskip}{6pt plus 2pt minus 1pt}}
}{% if KOMA class
  \KOMAoptions{parskip=half}}
\makeatother
\usepackage{color}
\usepackage{fancyvrb}
\newcommand{\VerbBar}{|}
\newcommand{\VERB}{\Verb[commandchars=\\\{\}]}
\DefineVerbatimEnvironment{Highlighting}{Verbatim}{commandchars=\\\{\}}
% Add ',fontsize=\small' for more characters per line
\usepackage{framed}
\definecolor{shadecolor}{RGB}{248,248,248}
\newenvironment{Shaded}{\begin{snugshade}}{\end{snugshade}}
\newcommand{\AlertTok}[1]{\textcolor[rgb]{0.94,0.16,0.16}{#1}}
\newcommand{\AnnotationTok}[1]{\textcolor[rgb]{0.56,0.35,0.01}{\textbf{\textit{#1}}}}
\newcommand{\AttributeTok}[1]{\textcolor[rgb]{0.13,0.29,0.53}{#1}}
\newcommand{\BaseNTok}[1]{\textcolor[rgb]{0.00,0.00,0.81}{#1}}
\newcommand{\BuiltInTok}[1]{#1}
\newcommand{\CharTok}[1]{\textcolor[rgb]{0.31,0.60,0.02}{#1}}
\newcommand{\CommentTok}[1]{\textcolor[rgb]{0.56,0.35,0.01}{\textit{#1}}}
\newcommand{\CommentVarTok}[1]{\textcolor[rgb]{0.56,0.35,0.01}{\textbf{\textit{#1}}}}
\newcommand{\ConstantTok}[1]{\textcolor[rgb]{0.56,0.35,0.01}{#1}}
\newcommand{\ControlFlowTok}[1]{\textcolor[rgb]{0.13,0.29,0.53}{\textbf{#1}}}
\newcommand{\DataTypeTok}[1]{\textcolor[rgb]{0.13,0.29,0.53}{#1}}
\newcommand{\DecValTok}[1]{\textcolor[rgb]{0.00,0.00,0.81}{#1}}
\newcommand{\DocumentationTok}[1]{\textcolor[rgb]{0.56,0.35,0.01}{\textbf{\textit{#1}}}}
\newcommand{\ErrorTok}[1]{\textcolor[rgb]{0.64,0.00,0.00}{\textbf{#1}}}
\newcommand{\ExtensionTok}[1]{#1}
\newcommand{\FloatTok}[1]{\textcolor[rgb]{0.00,0.00,0.81}{#1}}
\newcommand{\FunctionTok}[1]{\textcolor[rgb]{0.13,0.29,0.53}{\textbf{#1}}}
\newcommand{\ImportTok}[1]{#1}
\newcommand{\InformationTok}[1]{\textcolor[rgb]{0.56,0.35,0.01}{\textbf{\textit{#1}}}}
\newcommand{\KeywordTok}[1]{\textcolor[rgb]{0.13,0.29,0.53}{\textbf{#1}}}
\newcommand{\NormalTok}[1]{#1}
\newcommand{\OperatorTok}[1]{\textcolor[rgb]{0.81,0.36,0.00}{\textbf{#1}}}
\newcommand{\OtherTok}[1]{\textcolor[rgb]{0.56,0.35,0.01}{#1}}
\newcommand{\PreprocessorTok}[1]{\textcolor[rgb]{0.56,0.35,0.01}{\textit{#1}}}
\newcommand{\RegionMarkerTok}[1]{#1}
\newcommand{\SpecialCharTok}[1]{\textcolor[rgb]{0.81,0.36,0.00}{\textbf{#1}}}
\newcommand{\SpecialStringTok}[1]{\textcolor[rgb]{0.31,0.60,0.02}{#1}}
\newcommand{\StringTok}[1]{\textcolor[rgb]{0.31,0.60,0.02}{#1}}
\newcommand{\VariableTok}[1]{\textcolor[rgb]{0.00,0.00,0.00}{#1}}
\newcommand{\VerbatimStringTok}[1]{\textcolor[rgb]{0.31,0.60,0.02}{#1}}
\newcommand{\WarningTok}[1]{\textcolor[rgb]{0.56,0.35,0.01}{\textbf{\textit{#1}}}}
\usepackage{graphicx}
\makeatletter
\newsavebox\pandoc@box
\newcommand*\pandocbounded[1]{% scales image to fit in text height/width
  \sbox\pandoc@box{#1}%
  \Gscale@div\@tempa{\textheight}{\dimexpr\ht\pandoc@box+\dp\pandoc@box\relax}%
  \Gscale@div\@tempb{\linewidth}{\wd\pandoc@box}%
  \ifdim\@tempb\p@<\@tempa\p@\let\@tempa\@tempb\fi% select the smaller of both
  \ifdim\@tempa\p@<\p@\scalebox{\@tempa}{\usebox\pandoc@box}%
  \else\usebox{\pandoc@box}%
  \fi%
}
% Set default figure placement to htbp
\def\fps@figure{htbp}
\makeatother
\setlength{\emergencystretch}{3em} % prevent overfull lines
\providecommand{\tightlist}{%
  \setlength{\itemsep}{0pt}\setlength{\parskip}{0pt}}
\usepackage{booktabs}
\usepackage{longtable}
\usepackage{float}
\floatplacement{figure}{H}
\floatplacement{table}{H}
\usepackage{xcolor}
\definecolor{bawagblue}{HTML}{005493}
\definecolor{darkgreen}{HTML}{006400}
\definecolor{darkorange}{HTML}{FF8C00}
\definecolor{darkred}{HTML}{8B0000}
\usepackage{fancyhdr}
\pagestyle{fancy}
\fancyhf{}
\fancyhead[L]{\small BAWAG Group AG --- Risk Analysis 2025}
\fancyhead[R]{\small Group 4 | WU Wien 2026}
\fancyfoot[C]{\thepage}
\renewcommand{\headrulewidth}{0.4pt}
\usepackage{caption}
\captionsetup[table]{skip=4pt}
\AtBeginDocument{\hypersetup{colorlinks=true, linkcolor=bawagblue, urlcolor=bawagblue, citecolor=bawagblue}}
\usepackage{booktabs}
\usepackage{longtable}
\usepackage{float}
\usepackage{wrapfig}
\usepackage{threeparttable}
\usepackage{threeparttablex}
\usepackage{adjustbox}
\usepackage{array}
\usepackage{multirow}
\usepackage{tikz}
\usepackage{amsmath}
\usepackage{colortbl}
\usepackage{pdflscape}
\usepackage{tabu}
\usepackage[normalem]{ulem}
\usepackage{makecell}
\usepackage{xcolor}
\usepackage{bookmark}
\IfFileExists{xurl.sty}{\usepackage{xurl}}{} % add URL line breaks if available
\urlstyle{same}
\hypersetup{
  pdfauthor={Caterina Piacentini · Farkas Tallos · Balazs Limpek},
  hidelinks,
  pdfcreator={LaTeX via pandoc}}

\title{Credit Risk Measurement\\
Advanced Topics in Banking - Project 2\\
Group 4 \textbar{} WU Wien, SS 2026}
\author{Caterina Piacentini · Farkas Tallos · Balazs Limpek}
\date{2026-05-24}

\begin{document}
\maketitle

{
\setcounter{tocdepth}{3}
\tableofcontents
}
\section{Task 1 --- One-Year Migration Matrix
(Lando--Skodeberg)}\label{task-1-one-year-migration-matrix-landoskodeberg}

\emph{Author: Farkas Tallos}

\subsection{Case-Study Task}\label{case-study-task}

Based on the rating history of 10,000 representative borrowers,
determine the one-year migration matrix based on the Lando and Skodeberg
procedure. Determine the migration matrix using the generator matrix
based on at least the first five terms of the matrix exponential.

\subsection{Methodology}\label{methodology}

The Lando--Skodeberg {[}2002{]} estimator models rating dynamics as a
time-homogeneous Markov chain in continuous time. The generator matrix
\(\Lambda\) collects the instantaneous transition intensities. Its
maximum-likelihood estimator, given the full rating history of all
obligors, is

\[
\hat\lambda_{ij} \;=\; \frac{N_{ij}}{\int_0^T Y_i(s)\,ds}\,,\quad i \neq j,
\qquad
\hat\lambda_{ii} \;=\; -\sum_{j \neq i} \hat\lambda_{ij},
\]

where \(N_{ij}\) is the observed number of migrations from rating \(i\)
to rating \(j\) during the window \([0, T]\), and \(Y_i(s)\) is the
number of obligors in state \(i\) at time \(s\), so the denominator is
the total obligor-time spent in the from-state (slide p.126 of
\texttt{docs/01\_Slides\_QFin\_AT\ in\ Banking\_V1\_Handout.pdf}). The
one-year migration matrix is then obtained by exponentiating the
generator,

\[
P(\Delta t) \;=\; e^{\Lambda \cdot \Delta t} \;=\; \sum_{k=0}^{\infty}
\frac{(\Lambda \cdot \Delta t)^k}{k!},
\]

which we evaluate via a truncated Taylor expansion (slide p.132).
Default is modelled as an absorbing state, so the row of \(\Lambda\)
corresponding to \(\mathrm{D}\) is identically zero.

\subsection{State Space and Long-Format
Transitions}\label{state-space-and-long-format-transitions}

The 8-state rating universe used in this report is ordered from highest
quality to default, matching the lecture convention:

\begin{Shaded}
\begin{Highlighting}[]
\NormalTok{ratings\_order }\OtherTok{\textless{}{-}} \FunctionTok{c}\NormalTok{(}\StringTok{"AAA"}\NormalTok{, }\StringTok{"AA"}\NormalTok{, }\StringTok{"A"}\NormalTok{, }\StringTok{"BBB"}\NormalTok{, }\StringTok{"BB"}\NormalTok{, }\StringTok{"B"}\NormalTok{, }\StringTok{"CCC/C"}\NormalTok{, }\StringTok{"D"}\NormalTok{)}
\NormalTok{K }\OtherTok{\textless{}{-}} \FunctionTok{length}\NormalTok{(ratings\_order)}

\CommentTok{\# Wide {-}\textgreater{} long: one row per (Obligor, snapshot date, rating).}
\NormalTok{ratings\_long }\OtherTok{\textless{}{-}} \FunctionTok{melt}\NormalTok{(}
\NormalTok{    ratings,}
    \AttributeTok{id.vars         =} \StringTok{"Obligor"}\NormalTok{,}
    \AttributeTok{variable.name   =} \StringTok{"date"}\NormalTok{,}
    \AttributeTok{value.name      =} \StringTok{"rating"}\NormalTok{,}
    \AttributeTok{variable.factor =} \ConstantTok{FALSE}
\NormalTok{)}
\NormalTok{ratings\_long[, date   }\SpecialCharTok{:=} \FunctionTok{as.Date}\NormalTok{(date)]}
\NormalTok{ratings\_long[, rating }\SpecialCharTok{:=} \FunctionTok{factor}\NormalTok{(rating, }\AttributeTok{levels =}\NormalTok{ ratings\_order)]}
\FunctionTok{setorder}\NormalTok{(ratings\_long, Obligor, date)}

\CommentTok{\# Add the previous{-}month rating within each obligor; the first row of each}
\CommentTok{\# obligor receives NA because no prior snapshot exists.}
\NormalTok{ratings\_long[, prev\_rating }\SpecialCharTok{:=} \FunctionTok{shift}\NormalTok{(rating, }\AttributeTok{n =} \DecValTok{1}\DataTypeTok{L}\NormalTok{, }\AttributeTok{type =} \StringTok{"lag"}\NormalTok{),}
\NormalTok{             by }\OtherTok{=}\NormalTok{ Obligor]}
\NormalTok{transitions }\OtherTok{\textless{}{-}}\NormalTok{ ratings\_long[}\SpecialCharTok{!}\FunctionTok{is.na}\NormalTok{(prev\_rating)]}
\end{Highlighting}
\end{Shaded}

\begin{longtable}[t]{rlll}
\caption{\label{tab:task1-preview}Sample of obligor transitions (first five rows of the long table).}\\
\toprule
Obligor & date & prev\_rating & rating\\
\midrule
1 & 2025-01-31 & CCC/C & CCC/C\\
1 & 2025-02-28 & CCC/C & CCC/C\\
1 & 2025-03-31 & CCC/C & B\\
1 & 2025-04-30 & B & B\\
1 & 2025-05-31 & B & B\\
\bottomrule
\end{longtable}

Reshaping from wide to long yields one row per obligor-snapshot
observation, and lagging within \texttt{Obligor} produces the
\texttt{(prev\_rating\ -\textgreater{}\ rating)} pairs that feed the
migration count. Pairs where \texttt{prev\_rating} is \texttt{NA} (the
first observation of each obligor) carry no transition information and
are dropped. The data span 13 monthly snapshots from 2024-12-31 to
2025-12-31, giving 12 transitions per obligor and a total of
\(10{,}000 \times 12 = 120{,}000\) transition pairs.

\subsection{\texorpdfstring{Migration Counts
\(N_{ij}\)}{Migration Counts N\_\{ij\}}}\label{migration-counts-n_ij}

\begin{Shaded}
\begin{Highlighting}[]
\NormalTok{N }\OtherTok{\textless{}{-}} \FunctionTok{unclass}\NormalTok{(}\FunctionTok{table}\NormalTok{(}\AttributeTok{from =}\NormalTok{ transitions}\SpecialCharTok{$}\NormalTok{prev\_rating,}
                   \AttributeTok{to   =}\NormalTok{ transitions}\SpecialCharTok{$}\NormalTok{rating))}
\FunctionTok{storage.mode}\NormalTok{(N) }\OtherTok{\textless{}{-}} \StringTok{"double"}
\end{Highlighting}
\end{Shaded}

\begin{longtable}[t]{lrrrrrrrr}
\caption{\label{tab:task1-N-table}Migration count matrix $N_{ij}$: observed transitions from row-rating to column-rating.}\\
\toprule
 & AAA & AA & A & BBB & BB & B & CCC/C & D\\
\midrule
AAA & 5,860 & 37 & 4 & 1 & 1 & 0 & 0 & 0\\
AA & 7 & 11,417 & 85 & 11 & 0 & 0 & 0 & 0\\
A & 1 & 39 & 23,191 & 104 & 6 & 4 & 1 & 4\\
BBB & 1 & 11 & 120 & 35,982 & 164 & 28 & 7 & 10\\
BB & 2 & 0 & 5 & 102 & 23,440 & 163 & 22 & 19\\
\addlinespace
B & 0 & 0 & 1 & 2 & 52 & 11,927 & 62 & 65\\
CCC/C & 0 & 0 & 4 & 6 & 7 & 45 & 5,322 & 161\\
D & 0 & 0 & 0 & 0 & 0 & 0 & 0 & 1,497\\
\bottomrule
\end{longtable}

The count matrix is strongly diagonal-dominant: obligors are far more
likely to stay in their current rating than to migrate. Off-diagonal
mass concentrates immediately adjacent to the diagonal --- single-notch
moves dominate two-notch moves --- consistent with the rating-momentum
literature. The bottom row is identically zero off the diagonal because
default is absorbing in this state space, so no transitions originate
from \(\mathrm{D}\). The first few new defaults (4 from A, 10 from BBB,
19 from BB, 65 from B, 161 from CCC/C) are visible in the last column
and rise monotonically with risk, which is the qualitative behaviour the
IRB and simulation calibrations will inherit.

\subsection{\texorpdfstring{Time in State
\(Y_i\)}{Time in State Y\_i}}\label{time-in-state-y_i}

\begin{Shaded}
\begin{Highlighting}[]
\CommentTok{\# Add previous snapshot date alongside previous rating, then compute the}
\CommentTok{\# length of each interval (in years) for the from{-}state\textquotesingle{}s account.}
\NormalTok{ratings\_long[, prev\_date }\SpecialCharTok{:=} \FunctionTok{shift}\NormalTok{(date, }\AttributeTok{n =} \DecValTok{1}\DataTypeTok{L}\NormalTok{, }\AttributeTok{type =} \StringTok{"lag"}\NormalTok{),}
\NormalTok{             by }\OtherTok{=}\NormalTok{ Obligor]}
\NormalTok{ratings\_long[, dt }\SpecialCharTok{:=} \FunctionTok{as.numeric}\NormalTok{(date }\SpecialCharTok{{-}}\NormalTok{ prev\_date) }\SpecialCharTok{/} \FloatTok{365.25}\NormalTok{]}

\NormalTok{Y }\OtherTok{\textless{}{-}} \FunctionTok{tapply}\NormalTok{(ratings\_long}\SpecialCharTok{$}\NormalTok{dt, ratings\_long}\SpecialCharTok{$}\NormalTok{prev\_rating,}
\NormalTok{            sum, }\AttributeTok{na.rm =} \ConstantTok{TRUE}\NormalTok{)}
\NormalTok{Y }\OtherTok{\textless{}{-}}\NormalTok{ Y[ratings\_order]}
\end{Highlighting}
\end{Shaded}

\begin{longtable}[t]{lr}
\caption{\label{tab:task1-Y-table}Total time spent by all obligors in each rating $i$ (obligor-years).}\\
\toprule
Rating & Obligor-years\\
\midrule
AAA & 491.49\\
AA & 959.33\\
A & 1944.64\\
BBB & 3024.57\\
BB & 1977.84\\
\addlinespace
B & 1008.54\\
CCC/C & 461.51\\
D & 125.23\\
\bottomrule
\end{longtable}

Each obligor contributes the year-fraction between consecutive snapshots
to the rating it occupied at the start of that interval. Summing across
obligors gives the denominator of the Lando--Skodeberg estimator. The
exposure-time distribution mirrors the data's rating distribution: the
investment-grade buckets (BBB, A, BB) accumulate the most obligor-years,
while AAA and CCC/C are sparsely populated, which mechanically inflates
the variance of those rows of \(\Lambda\). The total
\(\sum_i Y_i \approx 9{,}993\) obligor-years almost exactly equals
10,000 obligors \(\times\) 1 year of window; the small shortfall
reflects the unequal calendar lengths of individual months.

\subsection{\texorpdfstring{Generator Matrix
\(\Lambda\)}{Generator Matrix \textbackslash Lambda}}\label{generator-matrix-lambda}

\begin{Shaded}
\begin{Highlighting}[]
\CommentTok{\# Off{-}diagonal entries: lambda\_\{ij\} = N\_\{ij\} / Y\_i (sweep divides row i by Y\_i).}
\NormalTok{Lambda }\OtherTok{\textless{}{-}} \FunctionTok{sweep}\NormalTok{(N, }\AttributeTok{MARGIN =} \DecValTok{1}\NormalTok{, }\AttributeTok{STATS =}\NormalTok{ Y, }\AttributeTok{FUN =} \StringTok{"/"}\NormalTok{)}

\CommentTok{\# Stays are not transitions: wipe the diagonal, then set it to the negative}
\CommentTok{\# row{-}sum so each row of Lambda sums to zero.}
\FunctionTok{diag}\NormalTok{(Lambda) }\OtherTok{\textless{}{-}} \DecValTok{0}
\FunctionTok{diag}\NormalTok{(Lambda) }\OtherTok{\textless{}{-}} \SpecialCharTok{{-}}\FunctionTok{rowSums}\NormalTok{(Lambda)}
\end{Highlighting}
\end{Shaded}

\begin{longtable}[t]{lrrrrrrrr}
\caption{\label{tab:task1-Lambda-table}Estimated generator $\Lambda$ (annualised transition intensities).}\\
\toprule
 & AAA & AA & A & BBB & BB & B & CCC/C & D\\
\midrule
AAA & -0.0875 & 0.0753 & 0.0081 & 0.0020 & 0.0020 & 0.0000 & 0.0000 & 0.0000\\
AA & 0.0073 & -0.1074 & 0.0886 & 0.0115 & 0.0000 & 0.0000 & 0.0000 & 0.0000\\
A & 0.0005 & 0.0201 & -0.0818 & 0.0535 & 0.0031 & 0.0021 & 0.0005 & 0.0021\\
BBB & 0.0003 & 0.0036 & 0.0397 & -0.1127 & 0.0542 & 0.0093 & 0.0023 & 0.0033\\
BB & 0.0010 & 0.0000 & 0.0025 & 0.0516 & -0.1583 & 0.0824 & 0.0111 & 0.0096\\
\addlinespace
B & 0.0000 & 0.0000 & 0.0010 & 0.0020 & 0.0516 & -0.1805 & 0.0615 & 0.0644\\
CCC/C & 0.0000 & 0.0000 & 0.0087 & 0.0130 & 0.0152 & 0.0975 & -0.4832 & 0.3489\\
D & 0.0000 & 0.0000 & 0.0000 & 0.0000 & 0.0000 & 0.0000 & 0.0000 & 0.0000\\
\bottomrule
\end{longtable}

Off-diagonal entries are instantaneous transition intensities, while the
diagonal equals the negative row-sum so mass is conserved. Row sums are
zero to machine precision (largest absolute deviation 1.1e-17).
Intensities grow sharply as we move down the rating ladder: the
magnitude of the CCC/C diagonal is roughly an order of magnitude larger
than the AAA diagonal, meaning CCC/C-rated obligors leave their bucket
much faster than high-grade names --- which, after exponentiation, will
translate into PDs differing by several orders of magnitude.

\subsection{Matrix Exponential --- One-Year Migration
Matrix}\label{matrix-exponential-one-year-migration-matrix}

We implement the truncated Taylor expansion of \(e^{\Lambda}\) with the
standard recursion
\(\mathrm{term}_k = \mathrm{term}_{k-1} \cdot \Lambda / k\), which
avoids any explicit factorial and uses one matrix multiplication per
iteration:

\begin{Shaded}
\begin{Highlighting}[]
\NormalTok{expm\_taylor }\OtherTok{\textless{}{-}} \ControlFlowTok{function}\NormalTok{(M, k\_max) \{}
\NormalTok{    K    }\OtherTok{\textless{}{-}} \FunctionTok{nrow}\NormalTok{(M)}
\NormalTok{    P    }\OtherTok{\textless{}{-}} \FunctionTok{diag}\NormalTok{(K)            }\CommentTok{\# k = 0 term (identity)}
\NormalTok{    term }\OtherTok{\textless{}{-}} \FunctionTok{diag}\NormalTok{(K)            }\CommentTok{\# running term  M\^{}k / k!}
    \ControlFlowTok{for}\NormalTok{ (k }\ControlFlowTok{in} \FunctionTok{seq\_len}\NormalTok{(k\_max)) \{}
\NormalTok{        term }\OtherTok{\textless{}{-}}\NormalTok{ (term }\SpecialCharTok{\%*\%}\NormalTok{ M) }\SpecialCharTok{/}\NormalTok{ k}
\NormalTok{        P    }\OtherTok{\textless{}{-}}\NormalTok{ P }\SpecialCharTok{+}\NormalTok{ term}
\NormalTok{    \}}
    \FunctionTok{dimnames}\NormalTok{(P) }\OtherTok{\textless{}{-}} \FunctionTok{dimnames}\NormalTok{(M)}
\NormalTok{    P}
\NormalTok{\}}

\NormalTok{P5  }\OtherTok{\textless{}{-}} \FunctionTok{expm\_taylor}\NormalTok{(Lambda, }\AttributeTok{k\_max =} \DecValTok{5}\NormalTok{)   }\CommentTok{\# case{-}study minimum}
\NormalTok{P20 }\OtherTok{\textless{}{-}} \FunctionTok{expm\_taylor}\NormalTok{(Lambda, }\AttributeTok{k\_max =} \DecValTok{20}\NormalTok{)  }\CommentTok{\# essentially exact}

\NormalTok{trunc\_error  }\OtherTok{\textless{}{-}} \FunctionTok{max}\NormalTok{(}\FunctionTok{abs}\NormalTok{(P5 }\SpecialCharTok{{-}}\NormalTok{ P20))}
\NormalTok{expm\_compare }\OtherTok{\textless{}{-}} \FunctionTok{max}\NormalTok{(}\FunctionTok{abs}\NormalTok{(P20 }\SpecialCharTok{{-}}\NormalTok{ expm}\SpecialCharTok{::}\FunctionTok{expm}\NormalTok{(Lambda)))}
\end{Highlighting}
\end{Shaded}

\begin{longtable}[t]{lrrrrrrrr}
\caption{\label{tab:task1-P5-table}One-year migration matrix $P^{(5)}$ from the 5-term Taylor expansion (\%).}\\
\toprule
 & AAA & AA & A & BBB & BB & B & CCC/C & D\\
\midrule
AAA & 91.648 & 6.840 & 1.057 & 0.253 & 0.188 & 0.010 & 0.002 & 0.003\\
AA & 0.665 & 89.927 & 8.090 & 1.244 & 0.044 & 0.015 & 0.003 & 0.011\\
A & 0.055 & 1.837 & 92.329 & 4.876 & 0.409 & 0.220 & 0.052 & 0.224\\
BBB & 0.035 & 0.363 & 3.626 & 89.562 & 4.771 & 1.007 & 0.223 & 0.412\\
BB & 0.090 & 0.015 & 0.325 & 4.531 & 85.674 & 7.035 & 1.011 & 1.319\\
\addlinespace
B & 0.002 & 0.002 & 0.119 & 0.321 & 4.403 & 83.898 & 4.459 & 6.796\\
CCC/C & 0.001 & 0.009 & 0.686 & 1.033 & 1.329 & 7.091 & 61.893 & 27.958\\
D & 0.000 & 0.000 & 0.000 & 0.000 & 0.000 & 0.000 & 0.000 & 100.000\\
\bottomrule
\end{longtable}

\begin{longtable}[t]{lrrrrrrrr}
\caption{\label{tab:task1-P20-table}One-year migration matrix $P^{(20)}$ from the 20-term Taylor expansion (\%).}\\
\toprule
 & AAA & AA & A & BBB & BB & B & CCC/C & D\\
\midrule
AAA & 91.648 & 6.840 & 1.057 & 0.253 & 0.188 & 0.010 & 0.002 & 0.003\\
AA & 0.665 & 89.927 & 8.090 & 1.244 & 0.044 & 0.015 & 0.003 & 0.011\\
A & 0.055 & 1.837 & 92.329 & 4.876 & 0.409 & 0.220 & 0.052 & 0.224\\
BBB & 0.035 & 0.363 & 3.626 & 89.562 & 4.771 & 1.007 & 0.223 & 0.412\\
BB & 0.090 & 0.015 & 0.325 & 4.531 & 85.674 & 7.035 & 1.011 & 1.319\\
\addlinespace
B & 0.002 & 0.002 & 0.119 & 0.321 & 4.403 & 83.898 & 4.459 & 6.796\\
CCC/C & 0.001 & 0.009 & 0.686 & 1.033 & 1.329 & 7.091 & 61.895 & 27.956\\
D & 0.000 & 0.000 & 0.000 & 0.000 & 0.000 & 0.000 & 0.000 & 100.000\\
\bottomrule
\end{longtable}

The case-study requires at least five terms of the series (slide p.132);
we additionally compute the 20-term expansion as a benchmark. The
maximum absolute element-wise difference between \(P^{(5)}\) and
\(P^{(20)}\) is 2e-05, i.e.~the 5-term truncation is already accurate to
about two parts in \(10^5\) on every cell. For external validation,
\(P^{(20)}\) agrees with \texttt{expm::expm(Lambda)} to 2.2e-16, which
is machine precision. The diagonal entries of \(P^{(20)}\) give the
one-year \emph{persistence} probability for each rating, falling from
\$\sim\$91.6\% for AAA to \$\sim\$61.9\% for CCC/C, and the last row
equals \((0, \dots, 0, 1)\) exactly, confirming that default remains
absorbing under the exponential map.

\subsection{One-Year PD Vector (with 5 bp
Floor)}\label{one-year-pd-vector-with-5-bp-floor}

\begin{Shaded}
\begin{Highlighting}[]
\NormalTok{PD\_raw }\OtherTok{\textless{}{-}}\NormalTok{ P20[, }\StringTok{"D"}\NormalTok{]}
\NormalTok{PD     }\OtherTok{\textless{}{-}} \FunctionTok{pmax}\NormalTok{(PD\_raw, }\FloatTok{0.0005}\NormalTok{)   }\CommentTok{\# 5 basis{-}point regulatory floor}
\end{Highlighting}
\end{Shaded}

\begin{longtable}[t]{lrrl}
\caption{\label{tab:task1-PD-table}One-year probabilities of default by rating, before and after the 5 bp regulatory floor.}\\
\toprule
Rating & PD\_raw & PD & Floored\\
\midrule
AAA & 0.000027 & 0.000500 & TRUE\\
AA & 0.000114 & 0.000500 & TRUE\\
A & 0.002237 & 0.002237 & FALSE\\
BBB & 0.004118 & 0.004118 & FALSE\\
BB & 0.013191 & 0.013191 & FALSE\\
\addlinespace
B & 0.067963 & 0.067963 & FALSE\\
CCC/C & 0.279564 & 0.279564 & FALSE\\
D & 1.000000 & 1.000000 & FALSE\\
\bottomrule
\end{longtable}

The raw AAA and AA PDs (approximately 2.7e-05 and 1.1e-04) lie below the
5 basis-point floor required by the regulator (CRR Art. 160(1)). This is
a well-known feature of investment-grade migration data: the
highest-quality buckets simply do not generate enough defaults in a
finite sample for the empirical estimator to be reliable. The floor is
the supervisor's way of acknowledging that even AAA names carry strictly
positive default risk. From A downward the empirical PD dominates the
floor by at least three orders of magnitude, with CCC/C reaching
\$\sim\$28\%. The full ordering
\(\text{PD}_{\text{AAA}} < \text{PD}_{\text{AA}} < \cdots <
\text{PD}_{\text{CCC/C}}\) is preserved, which is necessary for the
downstream regulatory and simulation models to behave sensibly.

\subsection{Summary --- Inputs for Downstream
Tasks}\label{summary-inputs-for-downstream-tasks}

The three objects carried forward into Tasks 3, 4 and 5 are:

\begin{itemize}
\tightlist
\item
  \textbf{\texttt{PD}} --- 8-vector of one-year PDs by rating, after the
  5 bp floor. Primary input to the IRB closed-form formula (Task 3) and
  to both Merton-style simulations (Tasks 4 and 5).
\item
  \textbf{\texttt{P20}} --- 8 \(\times\) 8 one-year migration matrix.
  Useful if a downstream task needs the full migration probabilities
  rather than just the default column.
\item
  \textbf{\texttt{Lambda}} --- the underlying generator. Useful for
  computing migration probabilities at any horizon \(\Delta t\) via
  \(P(\Delta t) = e^{\Lambda \Delta t}\).
\end{itemize}

\section{Task 2 --- Basel Standardized
Approach}\label{task-2-basel-standardized-approach}

\emph{Author: Farkas Tallos}

\subsection{Case-Study Task}\label{case-study-task-1}

Determine the credit risk for the portfolio of 100 borrowers based on
the Basel Standardized Approach. Proceed by using the existing ratings
as if they were external ratings.

\subsection{Regulatory Risk-Weight Lookup (CRR Art.
122)}\label{regulatory-risk-weight-lookup-crr-art.-122}

Under the Standardized Approach the risk weight applied to a corporate
exposure is a step function of the external credit assessment (CRR
Article 122; slide p.208 of
\texttt{docs/01\_Slides\_QFin\_AT\ in\ Banking\_V1\_Handout.pdf}).
Following the case-study instruction, we treat the internal ratings as
if they were ECAI ratings. The mapping from the eight-bucket rating
universe to risk weights, with defaulted exposures attracting the 150\%
weight applicable when specific provisions cover less than 20\% of the
outstanding amount, is:

\begin{Shaded}
\begin{Highlighting}[]
\NormalTok{rw\_table }\OtherTok{\textless{}{-}} \FunctionTok{data.table}\NormalTok{(}
    \AttributeTok{Rating =} \FunctionTok{c}\NormalTok{(}\StringTok{"AAA"}\NormalTok{, }\StringTok{"AA"}\NormalTok{,  }\StringTok{"A"}\NormalTok{,  }\StringTok{"BBB"}\NormalTok{, }\StringTok{"BB"}\NormalTok{,  }\StringTok{"B"}\NormalTok{,  }\StringTok{"CCC/C"}\NormalTok{, }\StringTok{"D"}\NormalTok{),}
    \AttributeTok{RW     =} \FunctionTok{c}\NormalTok{(}\FloatTok{0.20}\NormalTok{,  }\FloatTok{0.20}\NormalTok{,  }\FloatTok{0.50}\NormalTok{, }\FloatTok{1.00}\NormalTok{,  }\FloatTok{1.00}\NormalTok{,  }\FloatTok{1.50}\NormalTok{, }\FloatTok{1.50}\NormalTok{,    }\FloatTok{1.50}\NormalTok{)}
\NormalTok{)}
\end{Highlighting}
\end{Shaded}

\begin{longtable}[t]{lr}
\caption{\label{tab:task2-rw-table-render}CRR Art. 122 risk weights for corporate exposures.}\\
\toprule
Rating & Risk Weight\\
\midrule
AAA & 0.2\\
AA & 0.2\\
A & 0.5\\
BBB & 1.0\\
BB & 1.0\\
\addlinespace
B & 1.5\\
CCC/C & 1.5\\
D & 1.5\\
\bottomrule
\end{longtable}

\subsection{Per-Obligor Risk-Weighted
Assets}\label{per-obligor-risk-weighted-assets}

\begin{Shaded}
\begin{Highlighting}[]
\NormalTok{port\_sa }\OtherTok{\textless{}{-}} \FunctionTok{merge}\NormalTok{(portfolio, rw\_table, }\AttributeTok{by =} \StringTok{"Rating"}\NormalTok{, }\AttributeTok{all.x =} \ConstantTok{TRUE}\NormalTok{)}
\NormalTok{port\_sa[, RWA     }\SpecialCharTok{:=}\NormalTok{ Loan.Volume }\SpecialCharTok{*}\NormalTok{ RW]}
\NormalTok{port\_sa[, Capital }\SpecialCharTok{:=} \FloatTok{0.08} \SpecialCharTok{*}\NormalTok{ RWA]}
\FunctionTok{setorder}\NormalTok{(port\_sa, }\SpecialCharTok{{-}}\NormalTok{RWA)}
\end{Highlighting}
\end{Shaded}

\begin{longtable}[t]{rlrrrrr}
\caption{\label{tab:task2-top10}Ten largest obligors by RWA under the Standardized Approach.}\\
\toprule
Obligor & Rating & Industry & Loan.Volume & RW & RWA & Capital\\
\midrule
27 & CCC/C & 1 & 850,000 & 1.5 & 1,275,000 & 102,000\\
3 & BB & 2 & 850,000 & 1.0 & 850,000 & 68,000\\
21 & BB & 3 & 850,000 & 1.0 & 850,000 & 68,000\\
22 & BBB & 1 & 850,000 & 1.0 & 850,000 & 68,000\\
67 & BBB & 1 & 850,000 & 1.0 & 850,000 & 68,000\\
\addlinespace
13 & A & 3 & 850,000 & 0.5 & 425,000 & 34,000\\
23 & A & 1 & 850,000 & 0.5 & 425,000 & 34,000\\
93 & B & 1 & 184,000 & 1.5 & 276,000 & 22,080\\
5 & CCC/C & 1 & 174,000 & 1.5 & 261,000 & 20,880\\
96 & B & 3 & 163,000 & 1.5 & 244,500 & 19,560\\
\bottomrule
\end{longtable}

\subsection{Aggregated Breakdown by Rating
Bucket}\label{aggregated-breakdown-by-rating-bucket}

\begin{Shaded}
\begin{Highlighting}[]
\NormalTok{by\_rating }\OtherTok{\textless{}{-}}\NormalTok{ port\_sa[, .(}
    \AttributeTok{N       =}\NormalTok{ .N,}
    \AttributeTok{EAD     =} \FunctionTok{sum}\NormalTok{(Loan.Volume),}
    \AttributeTok{RW      =} \FunctionTok{unique}\NormalTok{(RW),}
    \AttributeTok{RWA     =} \FunctionTok{sum}\NormalTok{(RWA),}
    \AttributeTok{Capital =} \FunctionTok{sum}\NormalTok{(Capital)}
\NormalTok{), by }\OtherTok{=}\NormalTok{ Rating]}

\CommentTok{\# Order by the canonical rating order, not alphabetical.}
\NormalTok{by\_rating[, Rating }\SpecialCharTok{:=} \FunctionTok{factor}\NormalTok{(Rating, }\AttributeTok{levels =}\NormalTok{ ratings\_order)]}
\FunctionTok{setorder}\NormalTok{(by\_rating, Rating)}

\CommentTok{\# Append totals row.}
\NormalTok{totals }\OtherTok{\textless{}{-}} \FunctionTok{data.table}\NormalTok{(}
    \AttributeTok{Rating  =} \StringTok{"Total"}\NormalTok{,}
    \AttributeTok{N       =} \FunctionTok{sum}\NormalTok{(by\_rating}\SpecialCharTok{$}\NormalTok{N),}
    \AttributeTok{EAD     =} \FunctionTok{sum}\NormalTok{(by\_rating}\SpecialCharTok{$}\NormalTok{EAD),}
    \AttributeTok{RW      =} \FunctionTok{sum}\NormalTok{(by\_rating}\SpecialCharTok{$}\NormalTok{RWA) }\SpecialCharTok{/} \FunctionTok{sum}\NormalTok{(by\_rating}\SpecialCharTok{$}\NormalTok{EAD),}
    \AttributeTok{RWA     =} \FunctionTok{sum}\NormalTok{(by\_rating}\SpecialCharTok{$}\NormalTok{RWA),}
    \AttributeTok{Capital =} \FunctionTok{sum}\NormalTok{(by\_rating}\SpecialCharTok{$}\NormalTok{Capital)}
\NormalTok{)}
\NormalTok{by\_rating\_out }\OtherTok{\textless{}{-}} \FunctionTok{rbind}\NormalTok{(}
\NormalTok{    by\_rating[, .(}\AttributeTok{Rating =} \FunctionTok{as.character}\NormalTok{(Rating), N, EAD, RW, RWA, Capital)],}
\NormalTok{    totals}
\NormalTok{)}
\end{Highlighting}
\end{Shaded}

\begin{longtable}[t]{lrrrrr}
\caption{\label{tab:task2-by-rating-render}Standardized Approach breakdown by rating bucket (EUR, RW in decimal).}\\
\toprule
Rating & N & EAD & RW & RWA & Capital\\
\midrule
AAA & 3 & 279,000 & 0.20 & 55,800 & 4,464\\
AA & 8 & 1,410,000 & 0.20 & 282,000 & 22,560\\
A & 17 & 3,430,000 & 0.50 & 1,715,000 & 137,200\\
BBB & 30 & 4,821,000 & 1.00 & 4,821,000 & 385,680\\
BB & 27 & 4,428,000 & 1.00 & 4,428,000 & 354,240\\
\addlinespace
B & 9 & 887,000 & 1.50 & 1,330,500 & 106,440\\
CCC/C & 6 & 1,321,000 & 1.50 & 1,981,500 & 158,520\\
Total & 100 & 16,576,000 & 0.88 & 14,613,800 & 1,169,104\\
\bottomrule
\end{longtable}

The portfolio is concentrated in BBB (30 obligors, \$\sim\$29\% of EAD)
and BB (27 obligors, \$\sim\$27\% of EAD), both of which carry the
capital-neutral 100\% weight. The eight AA and three AAA names absorb
20\% weights, the seventeen A-rated obligors 50\%, while the fifteen
sub-investment-grade names (B and CCC/C combined) are penalised at
150\%. No obligors are in default in the current book, so the SA
exposure-weighted risk weight falls between the 75\% (everyone BBB) and
150\% (everyone CCC/C) extremes that would bracket this rating
distribution.

\subsection{Headline Capital Numbers}\label{headline-capital-numbers}

\begin{Shaded}
\begin{Highlighting}[]
\NormalTok{total\_EAD      }\OtherTok{\textless{}{-}} \FunctionTok{sum}\NormalTok{(port\_sa}\SpecialCharTok{$}\NormalTok{Loan.Volume)}
\NormalTok{total\_RWA      }\OtherTok{\textless{}{-}} \FunctionTok{sum}\NormalTok{(port\_sa}\SpecialCharTok{$}\NormalTok{RWA)}
\NormalTok{capital\_req\_SA }\OtherTok{\textless{}{-}} \FloatTok{0.08} \SpecialCharTok{*}\NormalTok{ total\_RWA}
\NormalTok{avg\_rw         }\OtherTok{\textless{}{-}}\NormalTok{ total\_RWA }\SpecialCharTok{/}\NormalTok{ total\_EAD}
\NormalTok{eff\_cap\_ratio  }\OtherTok{\textless{}{-}}\NormalTok{ capital\_req\_SA }\SpecialCharTok{/}\NormalTok{ total\_EAD}
\end{Highlighting}
\end{Shaded}

Total exposure at default is EUR 16,576,000, total RWA is EUR
14,613,800, and the minimum capital requirement (8\% of RWA) is EUR
1,169,104. The exposure-weighted average risk weight is 88.16\%, and the
effective capital-to-EAD ratio is 7.05\%.

\subsection{Interpretation and
Limitations}\label{interpretation-and-limitations}

The Standardized Approach delivers a minimum capital figure of roughly
EUR 1.17 m, or 7.05\% of total exposure. That number is mechanically
driven by the rating distribution and a coarse step-function of risk
weights; it ignores three economically important features of the
portfolio. \emph{First}, it makes no use of the empirical PDs we just
estimated in Task 1: a BBB and a BB obligor receive identical 100\%
weights despite their Lando--Skodeberg PDs differing by a factor of more
than three. \emph{Second}, it does not account for concentration risk;
the largest single exposure attracts exactly the same per-euro weight as
the smallest. \emph{Third}, it implicitly assumes independence --- there
is no correlation parameter and therefore no notion of tail or
unexpected loss. Tasks 3 (IRB), 4 (2-factor Merton with a common market
factor) and 5 (2-factor Merton with industry factors) progressively
relax these assumptions, and the resulting capital figures are compared
in Section \ref{sec:comp}.

\section{Task 3 --- Basel IRB Approach}\label{task-3-basel-irb-approach}

\emph{Author: Balazs Limpek}

\subsection{Case-Study Task}\label{case-study-task-2}

\subsection{Implementation}\label{implementation}

\subsection{Results}\label{results}

\section{Task 4 --- 2-Factor Merton with Common Market
Factor}\label{task-4-2-factor-merton-with-common-market-factor}

\emph{Author: Balazs Limpek}

\subsection{Case-Study Task}\label{case-study-task-3}

\subsection{Implementation}\label{implementation-1}

\subsection{Results}\label{results-1}

\section{Task 5 --- 2-Factor Merton with Industry
Factors}\label{task-5-2-factor-merton-with-industry-factors}

\emph{Author: Caterina Piacentini}

\subsection{Case-Study Task}\label{case-study-task-4}

\subsection{Implementation}\label{implementation-2}

\subsection{Results}\label{results-2}

\section{Comparison and Conclusions}\label{sec:comp}

\emph{Author: Caterina Piacentini}

\subsection{Capital Requirements Across
Methods}\label{capital-requirements-across-methods}

\subsection{Discussion}\label{discussion}

\subsection{Limitations and Caveats}\label{limitations-and-caveats}

\end{document}
